\section{Error analysis}\label{sec:error}

Throughout this section (IS3) holds; on Alfeld refinements with $k\ge3$ it is Lemma~\ref{st:IS3}.

The plan follows the two-dimensional analysis of \cite{PadhiLong2D}. Inserting the extended exact solution into the discrete equations leaves two residuals: a \emph{geometric} one, $\Gc$, of size $\hG^{3/2}$ in the energy dual, caused by the fact that $\ut$ does not vanish on $\Gh$ and that $\Gh$ is not $\Gamma$; and, for $\GS$ only, a \emph{pressure} residual $\Rc(v)=-\ip{\pt-\pbar}{v\cdot n}$, the traction that the printed method omits. The geometric residual gives the rate $\hG^{3/2}$. The pressure residual is of order one, and the question is how the discrete solution responds to it. The answer is the leak: the response is, to leading order, $(h/\mu)$ times a fixed divergence-free field $v_\varphi$ with normal trace $-(p-\pbar)$ on the wall (Lemma~\ref{er:vphi}). A cancellation that uses $\div v=0$ on each wall face (Lemma~\ref{er:cancel}) shows that this field has bounded energy, which is why the $H^1$ rate of $\GS$ is $1$ and not $\frac12$.

\subsection{Consistency}

\begin{lemma}[Error equation]\label{er:eq}
For $v\in\VR$,
\[
N_h(\ut,v)-(\pt,\div v)+\ip{\pt}{v\cdot n}=(\ft,v)+\Gc(v),\qquad
\Gc(v):=\ip{(\nabla\ut)^{\mathsf T}n}{v}-\ip{\ut}{\dn v}+\frac\mu h\ip{\ut}{v}-(F,v).
\]
\end{lemma}

\begin{proof}
Integrate $a_h(\ut,v)=(D(\ut),\nabla v)$ by parts on $\Oh$ ($v=0$ on $\Sigma$), use $\div D(\ut)=\Delta\ut$, $-\Delta\ut=\ft-F-\nabla\pt$, $-(\nabla\pt,v)=(\pt,\div v)-\ip{\pt}{v\cdot n}$, and $D(\ut)n-\dn\ut=(\nabla\ut)^{\mathsf T}n$.
\end{proof}

\begin{proposition}[The transposed-gradient term]\label{er:transposed}
On $F\in\FG$, $(\nabla\ut)^{\mathsf T}n_F=n(x^\ast)\bigl(\dn u(x^\ast)\cdot(n_F-n(x^\ast))\bigr)+O(\hG^2)=O(\hG)$. Hence, for $\gamma\ge\gamma_0$ and $v\in\VR$,
\[
|\ip{(\nabla\ut)^{\mathsf T}n}{v}|\le C\hG\norm v_{L^1(\Gh)}\le C\hG(h/\mu)^{1/2}\tnorm v=C\gamma^{-1/2}\hG^{3/2}\tnorm v .
\]
\end{proposition}

\begin{proof}
By Lemma~\ref{geom:wall}, $(\nabla u)^{\mathsf T}n_F=n\,(\dn u\cdot n_F)=n\,(\dn u\cdot(n_F-n))$ on $\Gamma$. Taylor-expand $\nabla\ut$ from $x^\ast$ to $x$ (distance $O(\hG^2)$) and use $|n_F-n(x^\ast)|\le C\hG$ (Theorem~\ref{geom:orient}). Then use $|\Gh|\le C$ and $\norm v^2_{\Gh}\le(h/\mu)\tnorm v^2$.
\end{proof}

In the $\norm{\nabla v}$-dual norm the transposed term is generically only of order $\hG$ in three dimensions: on the sphere, testing with the interpolant of $\chi\psi n$ gives $\sum_F|F|\psi(x_F^\ast)\,\dn u(x_F^\ast)\cdot(x_c-o_F)+O(\hG^2)$ by Lemma~\ref{geom:sag}(c), and the two-dimensional odd-symmetry cancellation is not available. The energy-dual bound above is all that the error analysis uses.

\begin{lemma}[Geometric consistency]\label{er:G}
For $\gamma\ge\gamma_0$: $|\Gc(v)|\le C(1+\gamma^{1/2})\hG^{3/2}\tnorm v$ for $v\in\VR$, and $|\Gc(v)|\le C\hG^{3/2}\norm{\nabla v}$ for $v\in\VO$.
\end{lemma}

\begin{proof}
The transposed term is Proposition~\ref{er:transposed}. By Lemmas~\ref{geom:ext} and~\ref{st:inverse}, $|\ip{\ut}{\dn v}|\le C\hG^2C_I(c_0\hG)^{-1/2}\norm{\nabla v}$; $\frac\mu h|\ip\ut v|\le(\mu/h)^{1/2}C\hG^2\tnorm v=C\gamma^{1/2}\hG^{3/2}\tnorm v$; and $|(F,v)|\le C\hG^2\norm{\nabla v}$. On $\VO$ the first and third terms vanish.
\end{proof}

\begin{lemma}[Fluxes]\label{er:flux}
For $v\in\Zh$, $\sum_i\int_{\Gamma_{h,i}}v\cdot n=0$, where $\Gamma_{h,i}:=\pi^{-1}(\Gamma_i)\cap\Gh$. If $m\ge2$, the map $v\mapsto(\int_{\Gamma_{h,i}}v\cdot n)_i$ is onto $\{x\in\R^m:\sum_ix_i=0\}$ for small $h$.
\end{lemma}

\begin{proof}
The first statement is $\int_{\partial\Oh}v\cdot n=\int_{\Oh}\div v=0$ with $v=0$ on $\Sigma$. For the second, apply Lemma~\ref{er:vphi} below with $q=\alpha_i$ constant on $\Gamma_i$, $\sum\alpha_i|\Gamma_i|=0$: its field has flux $-|\Gamma_i|\alpha_i+O(\hG^2+h^k)|\alpha|$ through $\Gamma_{h,i}$, by Lemma~\ref{geom:face}(c).
\end{proof}

\begin{corollary}[$\GS$ on $\Zh$]\label{er:GSZ}
For $v\in\Zh$ and every constant $c$, $N_h(\ut-u_h,v)=\Rc(v)+\Gc(v)$ with $\Rc(v):=-\ip{\pt-c}{v\cdot n}$.
\end{corollary}

\begin{remark}[The constant is global]\label{er:rem:global}
By Lemma~\ref{er:flux}, $\int_{\Gamma_{h,i}}v\cdot n$ is free on $\Zh$ for each component separately when $m\ge2$; only the total vanishes. So the constant in $\Rc$ can only be a single global one, and the relevant size of the wall pressure is $G=\norm{p-\pbar}_{L^2(\Gamma)}$ with the global mean $\pbar$, not the sum of per-component variances. Correspondingly, the correction in $\CNS$ uses the global mean $\pbG$; per-component means would make the method inconsistent. The two-dimensional statements and proofs are in \cite{PadhiLong2D}; they are dimension-free.
\end{remark}

\subsection{Approximation and the leak field}

Let
\[
\epsd:=\bigl(1+(\gamma\hG)^{1/2}\bigr)h^k+h^{k+1}\bigl(\hG^{-1/2}+(\mu/h)^{1/2}\bigr).
\]

\begin{lemma}[Approximation in $\Wh$]\label{er:approx}
$E_h:=\inf_{z\in\Wh}\tnorm{\ut-z}\le C\epsd$. For bounded $\gamma$ the second term is at most $Ch^k$ if $\hG\ge h^2$. If the outer data have zero net discrete flux, $\int_\Sigma\gI\cdot\nu=0$, it reduces to $C(1+\gamma^{1/2})\hG^{k+1/2}$. Every $z\in\Wh$ satisfies $\tnorm{\ut-z}\ge(\mu/h)^{1/2}|\int_\Sigma\gI\cdot\nu|/|\Gh|^{1/2}$.
\end{lemma}

\begin{proof}
Let $w_h:=I_h\ut$ be the Lagrange interpolant. Boundary tetrahedra have diameter at most $C\hG$, so $\norm{\nabla(\ut-w_h)}\le Ch^k$, $\norm{\ut-w_h}_{L^\infty(\Gh)}\le C\hG^{k+1}$ and $\norm{\dn(\ut-w_h)}_{L^\infty(\Gh)}\le C\hG^k$, whence $\tnorm{\ut-w_h}\le C(1+(\gamma\hG)^{1/2})h^k$ and $\norm{\div w_h}\le Ch^k$.
\emph{Net flux.} Since $\int_{\Gh}\ut\cdot n=-\int_\Sigma g\cdot\nu=0$, $m:=\int_{\Oh}\div w_h=\int_\Sigma(\gI-g)\cdot\nu+\int_{\Gh}(w_h-\ut)\cdot n$, and $|m|\le Ch^{k+1}$ by pointwise interpolation estimates. (Lagrange interpolation on triangles has no extra quadrature exactness in general, unlike the one-dimensional Newton--Cotes case of \cite{PadhiLong2D}.)
\emph{Bubble.} $b:=\sum_{F\in\FG}b_Fn_F$ satisfies $\int_{\Gh}b\cdot n=c_b|\Gh|$ with $c_b>0$, $\norm{\nabla b}+\norm{\div b}\le C\hG^{-1/2}$ (there are $O(\hG^{-2})$ bubbles, each with $H^1$-seminorm squared $O(\hG)$), $\norm b_{\Gh}\le C$ and $\sum_F\diam F\norm{\dn b}^2_F\le C\hG^{-1}$. With $a:=m/(c_b|\Gh|)$, $\tnorm{ab}\le C|a|(\hG^{-1/2}+(\mu/h)^{1/2})$.
\emph{Correction.} $q:=\div(w_h-ab)\in\Pih\cap L^2_0$ with $\norm q\le C\epsd$; (IS3) gives $v\in\VO$ with $\div v=q$ and $\tnorm v\le(1+C_I^2)^{1/2}\beta^{-1}\norm q$. Then $z:=w_h-ab-v\in\Wh$. The lower bound holds because $\int_{\Gh}z\cdot n=-\int_\Sigma\gI\cdot\nu$ for $z\in\Wh$, while $\int_{\Gh}\ut\cdot n=0$.
\end{proof}

\begin{lemma}[A divergence-free field with normal trace $q$]\label{er:w}
Let $q\in C^{k+1}(\Gamma)$ with $\int_\Gamma q=0$. There is $w\in H^{k+1}(\Oh\cup\Omega)^3$, of class $C^{k+1}$ on $\mathcal T_{r/16}$, with $\div w=0$, $w=0$ near $\Sigma$, and $w|_\Gamma=qN=-qn$.
\end{lemma}

\begin{proof}
Let $\bar q_i$ be the mean of $q$ on $\Gamma_i$ and $\phi_i\in C^{k+2,\alpha}(\Gamma_i)$ a solution of $\Delta_\Gamma\phi_i=q-\bar q_i$ (Schauder theory on the closed $C^{k+3}$ surface). Let $\chi\in C^\infty(\R)$ with $\chi=1$ on $(-\infty,r/8]$ and $\chi=0$ on $[r/4,\infty)$. On $\Gamma\times(-r,r)$ with volume $dA\,ds$ put $V:=(-\chi'(s)\nabla_\Gamma\phi_i,\ \chi(s)q)$, whose product divergence is $-\chi'\Delta_\Gamma\phi_i+\chi'q=\chi'\bar q_i$. Let $w_0$ be its Piola push-forward by $\Psi$, i.e.\ $\iota_{w_0}\mathrm{vol}=(\Psi^{-1})^\ast\iota_V(dA\wedge ds)$. Since $\Psi^\ast\mathrm{vol}=J\,dA\wedge ds$, $\div w_0=J^{-1}\chi'(s)\bar q_i$; at $s=0$, $D\Psi$ is the identity on $T\Gamma$ plus $\partial_s\mapsto N$, $J=1$ and $\chi'=0$, so $w_0=qN$ on $\Gamma$. Regularity: $\nabla_\Gamma\phi_i\in C^{k+1,\alpha}$ and $D\Psi,J\in C^{k+1}$, so $w_0\in C^{k+1}$.
$g:=\div w_0\in C^k$ is supported in $\{r/8\le d\le r/4\}\subset\Omega$, and $\int g=\sum_i\bar q_i|\Gamma_i|(\chi(r/4)-\chi(0))=-\int_\Gamma q=0$. Let $\mathcal W:=\{x\in\Omega:d(x)>r/16\}$. It is open and connected: the map $\Psi(y,s)\mapsto\Psi(y,\max(s,r/8))$ on $\Omega\cap\mathcal T_r$, extended by the identity, is a continuous map of the connected set $\Omega$ onto $\{x\in\Omega:d(x)\ge r/8\}$, and $\mathcal W$ is the union of this set with the connected collars $\Psi(\Gamma_i\times(r/16,r/8])$, each of which meets it. Join the $m$ connected sets $\Psi(\Gamma_i\times[r/8,r/4])$, whose union contains $\operatorname{supp}g$, by paths in $\mathcal W$; the union $K$ is compact and connected. Take $\psi\in C^\infty_c(\mathcal W)$ with $0\le\psi\le1$ and $\psi=1$ near $K$, a regular value $t\in(0,1)$ of $\psi$ (Sard), and let $U$ be the component of $\{\psi>t\}$ containing $K$: a bounded connected domain with $C^\infty$ boundary and $\overline U\subset\mathcal W$, so $\overline U$ has positive distance from $\Sigma$ and $d>r/16$ on $\overline U$. A Bogovski\u{\i} field $w_1\in H^{k+1}_0(U)^3$ with $\div w_1=-g$ exists \cite{CostabelMcIntosh2010,Galdi}. Put $w:=w_0+w_1$ (extended by zero); $w=w_0$ on $\mathcal T_{r/16}$. If $m=1$ or $q$ has zero mean on each component, $w_1=0$. %% CHECKED (v7): Costabel--McIntosh (arXiv:0808.2614 = Math. Z. 265) state in the introduction solvability of div u=f in W^{m,p}_0(Omega)^n for f in W^{m-1,p}_0, int f=0, on all bounded Lipschitz domains; their Cor. 3.4 / Lemma 4.3 / Thm 4.6 give H^s -> H^{s+1} for all s (notes/v7/cites/FINDINGS.txt [4]). No Galdi locator is cited here.
\end{proof}
For the model case and $q=(p-\pbar)|_\Gamma$ this is the explicit field $w=\chi(r)r^{-2}q(\omega)e_r-\chi'(r)r^{-1}\nabla_{S^2}\phi_0(\omega)$ in spherical coordinates, with $\Delta_{S^2}\phi_0=q$.

From now on $q:=(p-\pbar)|_\Gamma$, which requires $p|_\Gamma\in C^{k+1}$, and $w$ is the field of Lemma~\ref{er:w}.

\begin{lemma}[Discrete leak field]\label{er:vphi}
Let $b$ and $c_b$ be as in the proof of Lemma~\ref{er:approx}, and
\[
v_\varphi:=I_hw-ab-v_0,\qquad a:=\frac{\int_{\Gh}I_hw\cdot n}{c_b|\Gh|},
\]
with $v_0\in\VO$ given by (IS3) with $\div v_0=\div(I_hw-ab)$. Then $v_\varphi\in\Zh$ and
\begin{gather*}
\norm{v_\varphi-w}_{L^\infty(\Gh)}\le C\hG^{k+1},\qquad\norm{\nabla(v_\varphi-w)}\le Ch^k,\\
\norm{\dn(v_\varphi-w)}_{\Gh}\le Ch^k\hG^{-1/2},\qquad\sum_F\diam F\norm{\dn(v_\varphi-w)}^2_F\le Ch^{2k}.
\end{gather*}
\end{lemma}

\begin{proof}
Since $\div w=0$ on $\Oh$ and $w=0$ near $\Sigma$, $\int_{\Gh}w\cdot n=0$. So $|a|\le C\norm{I_hw-w}_{L^\infty(\Gh)}\le C\hG^{k+1}$ (boundary tetrahedra have diameter $\le C\hG$ and $w\in C^{k+1}$ there), and $\div(I_hw-ab)\in\Pih\cap L^2_0$ with norm at most $\sqrt3\norm{\nabla(I_hw-w)}+\sqrt3|a|\norm{\nabla b}\le Ch^k$. Hence $\norm{\nabla v_0}\le Ch^k$; as $v_0=0$ on $\Gh$, Lemma~\ref{st:inverse} gives $\norm{\dn v_0}_{\Gh}\le Ch^k\hG^{-1/2}$ and $\sum_F\diam F\norm{\dn v_0}^2_F\le Ch^{2k}$. The terms $I_hw-w$ and $ab$ contribute $O(\hG^{k+1})$ pointwise on $\Gh$, $O(h^k)$ in $H^1$ and $O(\hG^k)$ pointwise for $\dn$.
\end{proof}

These bounds replace the $W^{1,\infty}$ bound for the curl of a $C^1$ stream-function interpolant used in two dimensions \cite{PadhiLong2D}, in every use below.

\begin{lemma}[Normal-load cancellation]\label{er:cancel}
Let $S\in W^{1,\infty}$ on a neighbourhood of $\Gh$ with $|S-(S\cdot n_F)n_F|\le C\hG$ on every $F\in\FG$. (For instance $S=w$, since $w=qN+O(d)$ and $|n_F-n|\le C\hG$.) Then for $v\in\Zh$
\[
|\ip{S}{\dn v}|\le C\bigl(\norm v_{\Gh}+\hG^{1/2}\norm{\nabla v}\bigr)\le C\norm{\nabla v},
\]
with $C$ depending on $\norm S_{W^{1,\infty}}$, $k$ and the mesh constants.
\end{lemma}

\begin{proof}
\emph{Tangential part.} $|\ip{S-(S\cdot n_F)n_F}{\dn v}|\le C\hG\norm{\dn v}_{\Gh}\le C\hG^{1/2}\norm{\nabla v}$ by Lemma~\ref{st:inverse}.
\emph{Normal part.} On $F$, $v$ is a polynomial on $T_F$ with $\div v=0$, and $n_F$ is constant; so $\dn v\cdot n_F=-\div_Fv_F$, where $v_F$ is the part of $v$ tangential to $F$ and $\div_F$ the in-plane divergence. With $\sigma_F:=S\cdot n_F$ and $\nu_{\partial F}$ the in-plane outward conormal,
\[
\int_F\sigma_F\,\dn v\cdot n_F=\int_F\nabla_F\sigma_F\cdot v_F-\int_{\partial F}\sigma_F\,v\cdot\nu_{\partial F},
\]
and the first term is at most $C\norm v_{L^1(F)}$.
\emph{Edge terms.} Each edge $E$ of $\Gh$ lies in exactly two faces $F,F'$, and $v$ is continuous across $E$, so the two contributions combine to $\int_Ev\cdot(\sigma_F\nu_{E,F}+\sigma_{F'}\nu_{E,F'})$. With $t_E$ a unit tangent, $\nu_{E,F}=t_E\times n_F$ and $\nu_{E,F'}=-t_E\times n_{F'}$, so $\nu_{E,F}+\nu_{E,F'}=t_E\times(n_F-n_{F'})=O(\hG)$; and $\sigma_F-\sigma_{F'}=S\cdot(n_F-n_{F'})=O(\hG)$. So the coefficient is $O(\hG)$.
\emph{Summing.} $\norm v_{L^1(\partial F)}\le C(\diam F)^{-1}\norm v_{L^1(F)}$ for $v|_F\in P_k$, so $\sum_E\hG\norm v_{L^1(E)}\le C\norm v_{L^1(\Gh)}\le C\norm v_{\Gh}$. Finish with Lemma~\ref{geom:korn}.
\end{proof}

If boundary faces of $\Th$ subdivide the faces of $\Gh$ (for example on Worsey--Farin splits), the extra edges lie in a plane, where $n_F=n_{F'}$ and the conormals are opposite, so their contributions cancel exactly. On such meshes, however, (IS3) as formulated fails (Proposition~\ref{st:wf}), so the results of this section do not apply to them as they stand.

\subsection{The printed method}
We now quantify the leak. The energy norm contains the penalty-weighted wall term $(\mu/h)\norm{\cdot}_{\Gh}^2$, so a normal slip of size $(h/\mu)G$ costs energy of order $(h/\mu)^{1/2}G$; the $H^1$ norm does not see the wall term directly, and the slip costs only order $h/\mu$ there. Theorems~\ref{er:A}--\ref{er:B} make both statements precise, with matching lower bounds, and Corollary~\ref{er:Bp} identifies the leading term.

Throughout this subsection $\gamma\ge\gamma_0$. By Theorem~\ref{geom:orient}, Lemma~\ref{geom:face}(c) and $\pt(x)=p(x^\ast)+O(\hG^2)$ on $\Gh$,
\begin{equation}\label{er:riemann}
\norm{\pt-\pbar}_{\Gh}\le G(1+C\hG^2)+C\hG^2,\qquad \norm{v_\varphi}_{\Gh}=G(1+O(\hG^2))+O(\hG^{k+1}),\qquad \tnorm{v_\varphi}^2=\frac\mu h\norm{v_\varphi}^2_{\Gh}+O(1).
\end{equation}

\begin{theorem}[$\GS$: energy norm, exact rate $1/2$]\label{er:A}
$\tnorm{\ut-u_h}\le C\bigl[(\hG/\gamma)^{1/2}G+(1+\gamma^{1/2})\hG^{3/2}+\epsd\bigr]$. If $G>0$ and $h\le h_0(G)$, with $h_0$ independent of $\gamma\ge\gamma_0$,
\[
\tnorm{\ut-u_h}\ge(2M)^{-1}(\hG/\gamma)^{1/2}G-C(1+\gamma^{1/2})\hG^{3/2}.
\]
Here $(\hG/\gamma)^{1/2}=(h/\mu)^{1/2}$.
\end{theorem}

\begin{proof}
\emph{Upper bound.} For $z\in\Wh$, $e_h:=z-u_h\in\Zh$, and by Corollary~\ref{er:GSZ} with $c=\pbar$ and Lemma~\ref{st:coercive}, $c_1\tnorm{e_h}^2\le N_h(z-\ut,e_h)+\Rc(e_h)+\Gc(e_h)$. By \eqref{er:riemann}, $|\Rc(e_h)|\le(G+C\hG^2)(h/\mu)^{1/2}\tnorm{e_h}$; apply Lemmas~\ref{er:G} and~\ref{er:approx}.
\emph{Lower bound.} On $F$, $v_\varphi\cdot n_F=w(x^\ast)\cdot n_F+O(\hG^2)=-q(x^\ast)\,n(x^\ast)\cdot n_F+O(\hG^2)$ and $n(x^\ast)\cdot n_F=1+O(\hG^2)$; with Lemma~\ref{geom:face}(c), $\Rc(v_\varphi)=G^2+O(\hG^2+h^k)$. By continuity $M\tnorm{\ut-u_h}\tnorm{v_\varphi}\ge|\Rc(v_\varphi)|-|\Gc(v_\varphi)|$, and $|\Gc(v_\varphi)|\le C(1+\gamma^{1/2})\hG^{3/2}\tnorm{v_\varphi}$. By \eqref{er:riemann}, $|\Rc(v_\varphi)|/\tnorm{v_\varphi}\ge\frac12(h/\mu)^{1/2}G$ for $h\le h_0(G)$.
\end{proof}

\begin{theorem}[$\GS$: $H^1$ rate $1$]\label{er:C}
With $C_p:=C(1+\norm{p-\pbar}_{W^{1,\infty}(\Gamma)})$,
\[
\norm{\nabla(\ut-u_h)}\le C_p\,\frac h\mu+C\bigl[(1+\gamma^{1/2})\hG^{3/2}+\epsd\bigr].
\]
\end{theorem}

\begin{proof}
\emph{Split.} For $z\in\Wh$ write $z-u_h=\delta_R+\delta_G$ with $\delta_R,\delta_G\in\Zh$ solving $N_h(\delta_R,v)=\Rc(v)$ and $N_h(\delta_G,v)=\Gc(v)-N_h(\ut-z,v)$ for all $v\in\Zh$. Then $\norm{\nabla\delta_G}\le\tnorm{\delta_G}\le c_1^{-1}(C(1+\gamma^{1/2})\hG^{3/2}+ME_h)$.
\emph{Ansatz.} Put $r_h:=\delta_R-(h/\mu)v_\varphi\in\Zh$. Then $N_h(r_h,v)=m(v)+\ell(v)$ with
\[
m(v):=-\ip{(\pt-\pbar)n+v_\varphi}{v},\qquad\ell(v):=-\frac h\mu\bigl[a_h(v_\varphi,v)-\ip{\dn v_\varphi}v-\ip{v_\varphi}{\dn v}\bigr].
\]
\emph{Bound on $m$.} On $F$, $(\pt-\pbar)n_F+v_\varphi=q(x^\ast)(n_F-n(x^\ast))+O(\hG^2)=O(\hG)$, with nonzero face means in general (Lemma~\ref{geom:sag}(c)). So we use only the energy-dual bound $|m(v)|\le C_p\hG\norm v_{L^1(\Gh)}\le C_p\hG(h/\mu)^{1/2}\tnorm v$.
\emph{Bound on $\ell$.} $|a_h(v_\varphi,v)|\le C\norm{\nabla v}$; $|\ip{\dn v_\varphi}v|\le(C+Ch^k\hG^{-1/2})\norm v_{\Gh}\le(C_p+Ch^k\hG^{-1/2})\norm{\nabla v}$; and by Lemma~\ref{er:cancel} with $S=w$ and Lemma~\ref{er:vphi}, $|\ip{v_\varphi}{\dn v}|\le|\ip w{\dn v}|+C\hG^{k+1}\norm{\dn v}_{L^1(\Gh)}\le C_p\norm{\nabla v}$. So $|\ell(v)|\le[C_p(h/\mu)+C(h/\mu)h^k\hG^{-1/2}]\norm{\nabla v}$, and $(h/\mu)h^k\hG^{-1/2}=\gamma^{-1}h^k\hG^{1/2}\le Ch^k$.
\emph{Close.} Testing with $r_h$, $c_1\tnorm{r_h}\le C_p(\hG(h/\mu)^{1/2}+h/\mu)+Ch^k$, and $\hG(h/\mu)^{1/2}\le\frac12(\hG^2+h/\mu)$ gives $\tnorm{r_h}\le C_p(h/\mu+\hG^2)+C\epsd$. Hence $\norm{\nabla\delta_R}\le(h/\mu)\norm{\nabla v_\varphi}+\tnorm{r_h}\le C_p(h/\mu+\hG^2)+C\epsd$, and $\hG^2\le\hG^{3/2}$.
\end{proof}

The rate is $1$ and not $1/2$ because $\div v=0$ turns the normal part of the missing traction into a tangential derivative (Lemma~\ref{er:cancel}), as in two dimensions.

\begin{theorem}[$\GS$: boundary slip and $H^1$ lower bound]\label{er:B}
Let $G>0$ and $X:=C_p(\hG/\gamma)+(1+\gamma^{1/2})\hG^{3/2}+\epsd$. Suppose that, for a fixed small $\delta$,
\[
X\le\delta\,G\min(G,\hG^{1/2}),\qquad\hG/\gamma\le\delta G,\qquad\hG\le h_0(G).
\]
All three hold for fixed $\gamma$ and bounded $\rho$ as $\hG\to0$. Then
\begin{enumerate}[label=(\roman*)]
\item $c(\hG/\gamma)G\le\norm{\ut-u_h}_{L^2(\Gh)}\le C(\hG/\gamma)^{1/2}\tnorm{\ut-u_h}$;
\item $\norm{\nabla(\ut-u_h)}\ge c(\hG/\gamma)G-Ch^{k+1/2}$.
\end{enumerate}
\end{theorem}

\begin{proof}
Let $e_u:=\ut-u_h$ and $v:=v_\varphi$. Corollary~\ref{er:GSZ} gives
\[
\frac\mu h\ip{e_u}v=\Rc(v)+\Gc(v)-a_h(e_u,v)+\ip{\dn e_u}v+\ip{e_u}{\dn v}.
\]
(1) $\Rc(v)=G^2+O(\hG^2+h^k)$. (2) $|\Gc(v)|\le C(\hG+\gamma\hG^3+\hG^{3/2})$: the transposed term is $O(\hG)$ pointwise (Proposition~\ref{er:transposed}); in the penalty term, $\ut=d\,W(x^\ast)+O(d^2)$ is tangential to $\Gamma$ while $v=-qn(x^\ast)+O(\hG^2)$ is normal, so $\ut\cdot v=O(\hG^4)$. By the hypotheses, $|\Gc(v)|\le\delta G^2$. (3) $|a_h(e_u,v)|\le C\norm{\nabla e_u}\le CX\le C\delta G^2$ by Theorem~\ref{er:C}. (4) By Theorem~\ref{er:A}, $|\ip{e_u}{\dn v}|\le C_p\norm{e_u}_{\Gh}\le C_p(h/\mu)^{1/2}\tnorm{e_u}\le C_p[(h/\mu)G+X]\le2C_p\delta G^2$. (5) With Lemma~\ref{st:inverse}, $|\ip{\dn e_u}v|\le\bigl(\norm{\dn(\ut-z)}_{\Gh}+C\hG^{-1/2}\norm{\nabla(z-u_h)}\bigr)\norm v_{\Gh}\le C\hG^{-1/2}XG\le C\delta G^2$. For $\delta$ small, $|\ip{e_u}v|\ge\frac12(h/\mu)G^2$; divide by $\norm v_{\Gh}\le G+C\hG^2$. For (ii) let $E\in H^1(\Oh)^3$ lift $(g-\gI)|_\Sigma$, vanishing near $\Gh$, with $\norm E_{H^1}\le Ch^{k+1/2}$; then $\norm{e_u}_{\Gh}=\norm{e_u-E}_{\Gh}\le C_{\rm tr}(\norm{\nabla e_u}+\norm{\nabla E})$.
\end{proof}

\begin{corollary}[Leak law]\label{er:Bp}
Let $G>0$ and $\hG\to0$ with $h/\mu\to0$, $\gamma\hG\to0$ and $\epsd/(\hG/\gamma)^{1/2}\to0$. Then
\[
\ut-u_h=\frac h\mu\,v_\varphi+\vartheta_h,\qquad\tnorm{\vartheta_h}=o\bigl((h/\mu)^{1/2}\bigr),\qquad\norm{\vartheta_h}_{\Gh}=o(h/\mu),
\]
and consequently
\[
\frac{\norm{\ut-u_h}_{L^2(\Gh)}}{(h/\mu)G}\to1,\qquad\frac{\tnorm{\ut-u_h}}{(h/\mu)^{1/2}G}\to1 .
\]
To leading order the printed method replaces no-slip by the leak condition $(\mu/h)\,u_h\cdot n\approx p-\pbar$, with the global mean $\pbar$.
\end{corollary}

\begin{proof}
Write $\ut-u_h=(\ut-z)+\delta_G+\delta_R$ and $\delta_R=(h/\mu)v_\varphi+r_h$ as above, so $\vartheta_h=(\ut-z)+\delta_G+r_h$. Relative to $(h/\mu)^{1/2}G$: $\tnorm{r_h}\le C_p(h/\mu+\hG^2)+C\epsd$ contributes $O((h/\mu)^{1/2}+(\gamma\hG^3)^{1/2}+\epsd(\gamma/\hG)^{1/2})\to0$, and $\tnorm{\delta_G}+\tnorm{\ut-z}\le C[(1+\gamma^{1/2})\hG^{3/2}+\epsd]$ contributes $O((\gamma^{1/2}+\gamma)\hG+\epsd(\gamma/\hG)^{1/2})\to0$. The $\Gh$ norms carry an extra factor $(h/\mu)^{1/2}$. The limits follow from \eqref{er:riemann}.
\end{proof}

For fixed $\gamma$ and bounded $\rho$, Theorems~\ref{er:C} and~\ref{er:B} give $\norm{\nabla(\ut-u_h)}\asymp h/\mu$. Section~\ref{sec:leak} identifies the constant: $\norm{\nabla(\ut-u_h)}\mu/(hG)\to\norm{\nabla U}/G$, with $U$ the Stokes flow driven by the leak as Dirichlet data (Corollary~\ref{lk:gen}), and the leak limit holds for every $\gamma\ge\gamma_0$ (Theorem~\ref{lk:up}).

\subsection{The corrected method}
The corrected method removes the pressure residual. The only subtlety is the constant pressure mode: the pressure part of the traction must be applied with the wall mean removed, otherwise constants in $\Pih$ are invisible to the momentum equation.

\begin{proposition}[The plain traction is singular]\label{er:singular}
If $\ip{p_h-\pbG(p_h)}{v\cdot n}$ is replaced by $\ip{p_h}{v\cdot n}$, the square system ($\VR\times\Pih$ tested by $\VR\times\Pih$) has the kernel vector $(0,1)$ and is singular.
\end{proposition}

\begin{proof}
For $v\in\VR$, $-(1,\div v)+\ip1{v\cdot n}=-\int_{\partial\Oh}v\cdot n+\int_{\Gh}v\cdot n=0$, since $v=0$ on $\Sigma$.
\end{proof}
The $L^2_0$ variant, and the relation to \cite{FrachonNilssonZahedi2024}, are discussed in \cite{PadhiLong2D}; the arguments there are dimension-free.

With the mean removed, the discrete pressure appears in the momentum equation through the wall term, which is not controlled by the inf-sup condition alone. The proof therefore bounds the pressure by (IS3) in terms of the velocity error, and absorbs the wall term for penalties above a threshold $\gamma_2$ that depends on the inf-sup constant.

\begin{theorem}[$\CNS$: well posed, with Scott's rate]\label{er:D}
There is $\gamma_2\asymp\max(\gamma_0,\beta^{-2})$ such that for $\gamma\ge\gamma_2$ the method $\CNS$ has a unique solution and
\[
\tnorm{\ut-u_h}\le C\bigl[(1+\gamma^{1/2})\hG^{3/2}+\epsd\bigr].
\]
Consequently, for $\gamma\in[\gamma_2,\gamma_{\max}]$ and $\hG\ge h^2$,
\[
\norm{\ut-u_h}_{H^1(\Oh)}\le C\bigl(\hG^{3/2}+h^k\bigr),
\]
and with outer data of zero discrete net flux the proviso $\hG\ge h^2$ is not needed. On Alfeld refinements this holds for every $k\ge3$.
\end{theorem}

\begin{proof}
\emph{Consistency.} Let $B^\ast(q,v):=-(q,\div v)+\ip{q-\pbG(q)}{v\cdot n}$. For a constant $c$, $B^\ast(q+c,v)=B^\ast(q,v)-c\int_{\Gh}v\cdot n$. With $p^\ast:=\pt-\pbG(\pt)$, Lemma~\ref{er:eq} gives $N_h(\ut,v)+B^\ast(p^\ast,v)=(\ft,v)+\Gc(v)$, so with $e_u:=\ut-u_h$, $e_p:=p^\ast-p_h$,
\begin{equation}\label{er:cnserr}
N_h(e_u,v)+B^\ast(e_p,v)=\Gc(v)\qquad\forall v\in\VR .
\end{equation}
\emph{Split.} $e_p=\eta+\xi$ with $\eta:=p^\ast-\pi_hp^\ast$ ($\pi_h$ the $L^2$ projection onto $\Pih$) and $\xi\in\Pih$; $(\eta,\div v)=0$ for $v\in\VR$, and $\norm\eta_{L^\infty(\Gh)}\le C\hG^k$.
\emph{Pressure.} For $v\in\VO$, $(\xi-\bar\xi,\div v)=a_h(e_u,v)-\ip{e_u}{\dn v}-\Gc(v)$, where $\bar\xi$ is the mean over $\Oh$. With $|\ip{e_u}{\dn v}|\le(\rho/\mu)^{1/2}C_I\tnorm{e_u}\norm{\nabla v}$ and Lemma~\ref{er:G}, (IS3) gives $\beta\norm{\xi-\bar\xi}\le C(\tnorm{e_u}+\hG^{3/2})$.
\emph{Velocity.} For $z\in\Wh$, $e_h:=z-u_h\in\Zh$. Testing \eqref{er:cnserr} with $e_h$, and using $(e_p,\div e_h)=0$ and $\int_{\Gh}e_h\cdot n=0$,
\[
c_1\tnorm{e_h}^2\le ME_h\tnorm{e_h}+|\Gc(e_h)|+|\ip\eta{e_h\cdot n}|+|\ip{\xi-\bar\xi}{e_h\cdot n}|,
\]
with $|\ip\eta{e_h\cdot n}|\le C\hG^k(h/\mu)^{1/2}\tnorm{e_h}$ and, by Lemma~\ref{st:inverse} and $\rho/\mu\le(c_0\gamma)^{-1}$, $|\ip{\xi-\bar\xi}{e_h\cdot n}|\le C_I(c_0\gamma)^{-1/2}\norm{\xi-\bar\xi}\tnorm{e_h}$.
\emph{Absorb.} Insert the pressure bound and $\tnorm{e_u}\le E_h+\tnorm{e_h}$; for $\gamma\ge\gamma_2$ with $C_I(c_0\gamma_2)^{-1/2}C\beta^{-1}\le c_1/2$ the pressure term is absorbed.
\emph{Uniqueness.} For zero data, $z=0$, so $u_h=0$; then $p_h$ is constant, and \eqref{er:cnserr} reads $-c\int_{\Gh}v\cdot n=0$ for all $v$, so $c=0$ by the face bubble of Lemma~\ref{st:divrange}. The system is square.
\emph{$H^1$ bound.} For bounded $\gamma$ the flux term of $\epsd$ is $Ch^{k+1}\hG^{-1/2}\le Ch^k$ iff $\hG\ge h^2$; with zero net outer flux it is $C(1+\gamma^{1/2})\hG^{k+1/2}$ (Lemma~\ref{er:approx}). The $L^2$ part uses the lift $E$ of Theorem~\ref{er:B} and Lemma~\ref{geom:korn}. On Alfeld refinements (IS3) is Lemma~\ref{st:IS3}.
\end{proof}

\begin{remark}[Flux-corrected outer data]\label{er:rem:flux}
In the model case one may replace $\gI$ by $\tilde g_I:=\gI-\frac{m_R}{3|R|}x$ with $m_R:=\int_{\partial R}\gI\cdot\nu$; since $\int_{\partial R}x\cdot\nu=3|R|$, this has zero net flux. In two dimensions the corresponding $H^1$ bound holds uniformly in $\gamma$ \cite{PadhiLong2D}. In three dimensions Corollary~\ref{lk:gen} gives a $\gamma$-uniform $H^1$ bound for $\GS$ with $\gI$ itself (the flux term appears as $h^{k+1}\hG^{-1/2}$); for $\CNS$ we have not carried the uniformity over.
\end{remark}
